\chapter{Analysing Rules}
\label{ch:analysis}

RETE trades memory for speed: every partial match a rule can make is kept
in the network. Two rules that mean the same thing can differ several-fold
in the memory they use, depending on how their patterns are ordered and
how many facts each pattern lets through. This chapter describes the tools
\Morendo\ offers for comparing rules before the data arrives, and for
measuring the engine once it has.

\section{The topology cost}

\fn{topology-cost} assigns a rule a number that reflects the size of its
part of the network. The count is a static property of the rule and the
templates it uses; no facts are needed. It is computed as follows.

\begin{itemize}
\item Each pattern contributes 1 for its \emph{object type node} plus the
  cost of that node's successors: 1 for every alpha node and 4 for every
  join hanging directly off it. A template that many rules test in
  different ways therefore makes each of those rules look more expensive.
\item Each alpha node of the rule (a literal or predicate constraint)
  contributes 1.
\item Each join contributes 4, and a join on a not-equal binding
  (\fn{~?x}) contributes 5.
\end{itemize}

\begin{figure}[htb]
\centering
\includegraphics[width=0.85\textwidth]{duck_rete_cost_function_1}
\caption{Two rules over the \fn{duck} template compiled without sharing.
Each node has cost 1; the join-free rule \fn{middle\_age\_football\_fan}
costs 9 once its share of the type node's successors is counted.}
\label{fig:cost1}
\end{figure}

\begin{figure}[htb]
\centering
\includegraphics[width=0.85\textwidth]{duck_rete_sharing_1}
\caption{The same two rules with alpha-node sharing: \fn{gender is "m"}
and \fn{age > 50} are tested once for both.}
\label{fig:sharing1}
\end{figure}

Figures~\ref{fig:cost1} and \ref{fig:sharing1} (from \file{doc/}) show why
the number is only relative: the engine shares equal alpha nodes between
rules, so the node count of a rule depends on what else is loaded. Use the
cost to compare alternatives for the same rule in the same rule set, not
as an absolute.

\begin{lstlisting}[style=shell]
Morendo> (deftemplate person (slot name) (slot age))
true
Morendo> (defrule adult (person (name ?n) (age ?a&:(>= ?a 18)))
  => (printout t ?n crlf))
true
Morendo> (defrule pair (person (name ?n) (age ?a))
  (person (name ?m) (age ?b&:(> ?b ?a)))
  => (printout t ?m " older than " ?n crlf))
true
Morendo> (topology-cost adult)
Morendo> (topology-cost pair)
Morendo> (rules)
adult "" salience:100 version: no-agenda:false temporal-activation:false cost-value:10
pair "" salience:100 version: no-agenda:false temporal-activation:false cost-value:24
for a total of 2
Morendo> (average-cost)
17
\end{lstlisting}

\fn{pair} joins the \fn{person} template with itself: two object
conditions, a join and a predicate that must be evaluated on every
combination. \fn{adult} tests one alpha constraint. The
\fn{cost-value} column of \fn{(rules)} and the \fn{;; topology-cost}
line printed by \fn{ppdefrule} show the last value computed;
\fn{topology-cost-all} computes it for every rule and
\fn{average-cost} averages the current module. \fn{topology-cost}
accepts several rule names at once.

\fn{partial-match-cost} estimates the run-time side: how many partial
matches a rule will keep, given the number of distinct values each slot
has. Tell the engine those counts first with
\fn{(set-distinct-count template slot n)}.

\section{Generating test facts}

\fn{(generate-facts rule)} constructs a set of facts that satisfies a
rule's patterns, from the literal constraints in the rule and default
values for the rest; \fn{(test-rule rule)} generates them, asserts them and
fires, which is a quick way to check that a freshly written rule can
activate at all.

\section{Profiling the engine}

Profiling counts and times the engine's own operations. Switch it on
before loading data:

\begin{lstlisting}[style=shell]
Morendo> (profile all)
Morendo> (assert (person (name "ann") (age 34)))
2
Morendo> (assert (person (name "bob") (age 40)))
3
Morendo> (fire)
bob older than ann
bob
ann
3
Morendo> (print-profile)
fire ET=0 ms
assert ET=6 ms
retract ET=0 ms
add Activation ET=0 ms
remove Activation ET=0 ms
Activation added=3
Activation removed=0
Average cube query=0 ms
\end{lstlisting}

\fn{profile} and \fn{unprofile} accept \fn{all}, \fn{assert-fact},
\fn{retract-fact}, \fn{fire}, \fn{add-activation} and
\fn{remove-activation}, and several may be given at once. Cubes have their
own profiling functions (Chapter~\ref{ch:molap}). The counters belong to
the engine: two engines profiling in one JVM keep separate numbers, and
from Java they are read through \fn{engine.getProfileStats()}.
\fn{(reset)} and \fn{(clear)} zero them.

\section{Looking inside the network}

\fn{(rule-matches name)} prints, node by node, the facts a rule's alpha
and join nodes currently hold; \fn{(right-matches)} prints the right-hand
memories of every join, and \fn{(matches)} every memory in the network.
These are the tools for answering ``why did this rule not fire'': if a
fact is missing from an alpha memory it failed a constraint, and if it is
in both inputs of a join but not in the join's output the bindings did not
agree.

\fn{(mem-used)}, \fn{(mem-free)} and \fn{(mem-total)} print the JVM's
memory figures, and \fn{(gc)} asks for a collection; together with
\fn{(fact-count)} they give a rough memory cost per fact when loading a
large data set.
